\documentclass[10pt,a4paper]{amsart}
\usepackage[margin=19mm]{geometry}
\usepackage{amsmath,amssymb,mathtools}
\usepackage[scr=rsfs,cal=euler]{mathalfa}
\usepackage{xfrac}
\usepackage[unicode,hidelinks,bookmarks=false]{hyperref}
\newcommand{\ie}{i.e.\ }
\newcommand{\cf}{cf.\ }
\newcommand{\resp}{respectively\ }
\newcommand{\Subsection}{\subsection}
\providecommand{\nobreakdash}{\nobreak\hbox{-}}
\setlength{\emergencystretch}{2em}
\multlinegap0pt
\usepackage{bm}
\usepackage{bbm}
\usepackage{dsfont}
\usepackage{xspace}
\usepackage{cancel}
\usepackage{mathtools}
\usepackage{amsthm}
\makeatletter
\@ifundefined{theorem}{\newtheorem{theorem}{Theorem}}{}
\@ifundefined{lemma}{\newtheorem{lemma}[theorem]{Lemma}}{}
\makeatother
\title[Bank-Laine functions with preassigned number of zeros]{Bank-Laine functions with preassigned number of zeros}
\author{Yueyang Zhang}
\date{Source: arXiv:2311.13618v4 (CC BY 4.0)}
\begin{document}
\maketitle

\noindent\emph{Source and licence.} Bank-Laine functions with preassigned number of zeros. Version arXiv:2311.13618v4 (CC BY 4.0), distributed under the Creative Commons Attribution 4.0 International licence, as recorded in the arXiv licence metadata for this exact version and shown on the abstract page https://arxiv.org/abs/2311.13618v4. Full rights notice: licenses/arxiv-2311.13618-CC-BY-4.0.txt.\par\smallskip

\noindent\textit{Editorial scope.} Counted unit: the Lemma whose source label is \texttt{Lemma 0} in the article (source0.tex lines 280--293, source label \texttt{Lemma 0}), delivered together with the article's own complete original proof of it (source0.tex lines 295--372, one \texttt{proof} environment of 78 lines). No mathematics is written, repaired or re-derived: statement and proof are the article's own text line for line. Required context retained uncounted: recuree1pm9 (lines 161--165); recuree1pm10 (lines 167--171). Every definition, display and label of those lines is retained unchanged. Omitted from this package and not redistributed: 1--160, 166--166, 172--279, 294--294, 373--2384. Nothing inside a retained excerpt is omitted or altered: every retained body is delivered exactly as the article's lines are written, so no omission receipt of any kind accompanies this package. Citations: the retained excerpts carry no citation command at all, so no bibliography entry of the article is transcribed and no reference is redistributed. Reference adaptation: none was needed and none was made; every reference inside the delivered unit resolves inside it, so no unresolved reference or citation is produced. Printed slips kept literal: no printed word, symbol, hypothesis or number of the counted statement or of its proof was altered. Layout and class: the article's own class is \texttt{amsart} with packages including amsfonts, amssymb, amsmath, amsthm, cite; the wrapper is the lane's pinned \texttt{amsart} editorial wrapper at 10pt on A4, which restates the theorem-like environments the retained text needs and adds the article's own macro definitions that the retained text needs (none), with extra wrapper packages (bm, bbm, dsfont, xspace, cancel, mathtools). The delivered numbering is the unit's own. Assets: no retained span contains a figure environment, a graphics-inclusion command, a PDF-page inclusion, a table or a TikZ picture; no image file is copied into this package, the package declares no asset dependency and no image file is redistributed with it. Licence: version arXiv:2311.13618v4 of this article is distributed under the Creative Commons Attribution 4.0 International licence, as recorded in the arXiv licence metadata for this exact version and shown on the abstract page https://arxiv.org/abs/2311.13618v4. A full rights notice is delivered at licenses/arxiv-2311.13618-CC-BY-4.0.txt. Characters: every retained excerpt is pure ASCII, so the delivered engine, XeLaTeX with its default Latin Modern text font, reports no missing character.
\begin{equation}\label{recuree1pm9}
\begin{split}
g_{m,n}(z)=h_{m,n}(e^z)=\left(\frac{\sum_{j=0}^{2n}B_je^{jz}}{\sum_{i=0}^{m}A_ie^{iz}}\right)e^{e^{z}}
\end{split}
\end{equation}

\begin{equation}\label{recuree1pm10}
\begin{split}
g'_{m,n}(z)=\frac{1}{\binom{m+2n}{m}(m+2n)!}\frac{e^{e^{z}+(m+2n+1)z}}{\left(\sum_{i=0}^{m}A_ie^{iz}\right)^2}.
\end{split}
\end{equation}

\begin{lemma}\label{Lemma 0}
The asymptotic behaviors of the diffeomorphism $\phi$ satisfy:
\begin{eqnarray}
&\phi(x)=x+O(e^{-x/2}),        \quad\quad\quad\quad\quad &\phi'(x)=1+O(e^{-x/2}), \quad\quad x\to \infty,  \label{recuree1pm11}\\
&\phi(x)=\kappa x+c+O(e^{-\delta|x|/2}),    \quad\quad &\phi'(x)=\kappa+O(e^{-\delta|x|/2}),  \quad x\to -\infty,             \label{recuree1pm12}
\end{eqnarray}
with
\begin{equation*}%\label{recuree1pm12fu}
\begin{split}
c=\frac{1}{N}\log \left[\frac{\binom{m_1+2n_1}{m_1}}{\binom{m_2+2n_2}{m_2}}\frac{N!}{M!}\right] \quad \text{and} \quad \delta=\frac{1}{2}\min\{1,\kappa\}.
\end{split}
\end{equation*}

\end{lemma}

\begin{proof}
In order to prove \eqref{recuree1pm11}, we note that
\begin{equation*}%\label{recuree1pm13}
\begin{split}
\log g_{m_2,n_2}(x)=e^{x}+O(x)=e^{x}\left[1+O(xe^{-x})\right], \quad x\to \infty.
\end{split}
\end{equation*}
The equation $g_{m_2,n_2}(x)=(g_{m_1,n_1}\circ\phi)(x)$ easily implies that $2x/3\leq \phi(x)\leq 2x$ for large $x$. Thus we also have
\begin{equation*}%\label{recuree1pm14}
\begin{split}
\log g_{m_1,n_1}(\phi(x))&=e^{\phi(x)}\left[1+O(\phi(x)e^{-\phi(x)})\right]=e^{\phi(x)}\left[1+O(\phi(x)e^{-2x/3})\right], \quad x\to \infty.
\end{split}
\end{equation*}
Combining the last two equations we obtain
\begin{equation*}%\label{recuree1pm15}
\begin{split}
e^{\phi(x)-x}=1+O(xe^{-2x/3}), \quad x\to \infty,
\end{split}
\end{equation*}
from which the first statement in \eqref{recuree1pm11} easily follows. For the second statement in \eqref{recuree1pm11} we use
\begin{equation}\label{recuree1pm16  bef0}
\begin{split}
\phi'=\frac{g'_{m_2,n_2}}{g_{m_2,n_2}}\frac{g_{m_1,n_1}\circ\phi}{g'_{m_1,n_1}\circ\phi}=e^{Mx-N\phi(x)}\frac{A_{m_2}B_{2n_2}}{A_{m_1}B_{2n_1}}\frac{\sum_{i=0}^{m_1}A_ie^{i\phi(x)}}{\sum_{i=0}^{m_2}A_ie^{ix}}\frac{\sum_{j=0}^{2n_1}B_je^{j\phi(x)}}{\sum_{j=0}^{2n_2}B_je^{jx}},
\end{split}
\end{equation}
so that, together with the first statement in \eqref{recuree1pm11},
\begin{equation}\label{recuree1pm17 bef1}
\begin{split}
\phi'&=e^{Mx-N\phi(x)+(m_1+2n_1)\phi(x)-(m_2+2n_2)x}\frac{1+O(e^{-\phi(x)})}{1+O(e^{-x})}\frac{1+O(e^{-\phi(x)})}{1+O(e^{-x})}\\
&=e^{O(e^{-x/2})}\left[1+O(e^{-\phi(x)})+O(e^{-x})\right]=1+O(e^{-x/2}), \quad x\to \infty.
\end{split}
\end{equation}
In order to prove \eqref{recuree1pm12} we note from \eqref{recuree1pm9} and \eqref{recuree1pm10} that
\begin{equation*}%\label{recuree1pm18}
\begin{split}
h'_{m,n}(w)=\frac{1}{\binom{m+2n}{m}(m+2n)!}\frac{e^ww^{m+2n}}{Q_m(w)^2}.
\end{split}
\end{equation*}
Since $h_{m,n}(w)$ is analytic on $[0,\infty)$, then by taking the derivatives of $h'_{m,n}(w)$ successively we find that $h'_{m,n}(0)=\cdots=h^{(m+2n)}_{m,n}(0)=0$. Thus by Peano's version of Taylor's theorem we have
\begin{equation*}%\label{recuree1pm19}
\begin{split}
h_{m,n}(w)=R_{m,n}(w)e^{w}=1+\frac{w^{m+2n+1}}{\binom{m+2n}{m}(m+2n+1)!}+O(w^{m+2n+2}), \quad w\to 0.
\end{split}
\end{equation*}
Hence
\begin{equation*}%\label{recuree1pm20}
\begin{split}
g_{m,n}(x)&=1+\frac{e^{(m+2n+1)x}}{\binom{m+2n}{m}(m+2n+1)!}+O(e^{(m+2n+2)x})\\
&=1+\frac{e^{(m+2n+1)x}}{\binom{m+2n}{m}(m+2n+1)!}\left[1+O(e^{x})\right], \quad x\to-\infty.
\end{split}
\end{equation*}
The equation $g_{m_2,n_2}(x)=g_{m_1,n_1}(\phi(x))$ now yields
\begin{equation*}%\label{recuree1pm21}
\begin{split}
\frac{\binom{m_2+2n_2}{m_2}}{\binom{m_1+2n_1}{m_1}}\frac{M!}{N!}e^{N\phi(x)-Mx}=1+O(e^{x})+O(e^{\phi(x)}), \quad x\to-\infty
\end{split}
\end{equation*}
and hence
\begin{equation*}%\label{recuree1pm22}
\begin{split}
\phi(x)=\frac{M}{N}x+\frac{1}{N}\log \left[\frac{\binom{m_1+2n_1}{m_1}}{\binom{m_2+2n_2}{m_2}}\frac{N!}{M!}\right]+O(e^{x})+O(e^{\phi(x)}), \quad x\to-\infty,
\end{split}
\end{equation*}
which gives the first statement in \eqref{recuree1pm12}. For the second statement in \eqref{recuree1pm12}, recall that $A_{m}B_{2n}=\frac{m!(2n)!}{[(m+2n)!]^2}=\frac{1}{\binom{m+2n}{m}(m+2n)!}$. Similarly as for the estimates in \eqref{recuree1pm17 bef1}, \eqref{recuree1pm16  bef0} now takes the form
\begin{equation*}%\label{recuree1pm23}
\begin{split}
\phi'(x)=\frac{A_{m_2}B_{2n_2}}{A_{m_1}B_{2n_1}}e^{Mx-N\phi(x)}\frac{1+O(e^{\phi(x)})}{1+O(e^{x})}\frac{1+O(e^{\phi(x)})}{1+O(e^{x})}
\end{split}
\end{equation*}
and thus, by the first statement in \eqref{recuree1pm12}, we obtain
\begin{equation*}%\label{recuree1pm23}
\begin{split}
\phi'(x)&=\frac{A_{m_2}B_{2n_2}}{A_{m_1}B_{2n_1}}e^{Mx-N(\kappa x+c+O(e^{-\delta|x|/2}))}\left[1+O(e^{\phi(x)})+O(e^{x})\right]\\
&=\frac{A_{m_2}B_{2n_2}}{A_{m_1}B_{2n_1}}e^{-Nc+O(e^{-\delta|x|/2})}\left[1+O(e^{\phi(x)})+O(e^{x})\right]=\kappa+O(e^{-\delta|x|/2}).
\end{split}
\end{equation*}

\end{proof}

\end{document}
